% 10 slides, 10 minutes (spec §12). Built by `make report` with the figures and the
% result macros of the paper: nothing on these slides is typed by hand.
\documentclass[aspectratio=169,11pt]{beamer}
\usetheme{default}
\usecolortheme{seahorse}
\setbeamertemplate{navigation symbols}{}
\setbeamertemplate{footline}[frame number]
\usepackage[T1]{fontenc}
\usepackage{lmodern,booktabs,graphicx,amsmath}
\graphicspath{{figures/}}
\newcommand{\Hz}{H_0}
\newcommand{\fig}[2]{\IfFileExists{figures/#1}{\includegraphics[height=#2\textheight]{#1}}{\fbox{\parbox{0.6\linewidth}{\centering\small figure \texttt{#1}: built by \texttt{make report}}}}}
\IfFileExists{tables/macros.tex}{\input{tables/macros.tex}}{\newcommand{\resultHOne}{pending}\newcommand{\resultHTwo}{pending}\newcommand{\resultHThree}{pending}}
\providecommand{\resultFOne}{pending}
\providecommand{\resultFTwo}{pending}

\title{One exponent to rule them all?}
\subtitle{An asset-by-asset test of the scaling laws of order flow, volume, volatility and
impact on crypto markets}
\author{nidhal.V}
\date{Data: Binance Vision, CC BY-NC-SA 4.0 --- not affiliated with Binance}

\begin{document}

\begin{frame}
  \titlepage
\end{frame}

\begin{frame}{The question}
  A 2026 theory (Muhle-Karbe, Ouazzani Chahdi, Rosenbaum \& Szymanski): one number,
  the persistence $\Hz$ of the signed flow of core orders, fixes three other exponents.
  \medskip
  \begin{center}
  \begin{tabular}{lcc}
    \toprule
    Quantity & Predicted exponent & if $\Hz=3/4$ \\
    \midrule
    Signed flow (Hurst) & $\Hz$ & 0.75 \\
    Volume rate (Hurst) & $\Hz-\tfrac12$ & 0.25 \\
    Extraday volatility (Hurst) & $2\Hz-\tfrac32$ & 0 \\
    Impact of a meta-order (time exponent) & $2-2\Hz$ & 0.5 \\
    \bottomrule
  \end{tabular}
  \end{center}
  \medskip
  Tested so far only \alert{on average}, over 40 stocks. Does it hold \alert{asset by asset}?
\end{frame}

\begin{frame}{What is new here}
  \begin{itemize}
    \item The four exponents estimated on the \alert{same data}, for each of 60 crypto pairs,
      spot and perpetual.
    \item A \alert{joint} cross-sectional test of slopes and intercepts, corrected for the
      errors in $\widehat{\Hz}$ (Deming regression, synchronized block bootstrap, Wald, Holm).
    \item A different market: open 24 hours, the side of every trade reported by the exchange.
    \item Estimators validated \alert{by simulation} with microstructure noise before touching
      real data; hypotheses and decision rule \alert{pre-registered} (git tag
      \texttt{prereg-v1}); confirmation sample frozen (\texttt{config-frozen}).
  \end{itemize}
\end{frame}

\begin{frame}{Data}
  \begin{itemize}
    \item Binance Vision aggTrades, 60 USDT pairs ranked by median daily notional, spot and
      USD-M perpetuals, November 2024 -- October 2026 (12 months development, 12 months
      confirmation); up to 5 years of 1-minute klines for daily volatility.
    \item Every file: SHA-256, schema, id continuity, prices, OHLC against klines, jumps;
      every pair-day ends reduced, blocked or explained (\texttt{ofs quality}).
    \item Reduced to 1-s and 10-s windows and additive daily sufficient statistics, then
      deleted: the bootstrap never reads a trade again.
  \end{itemize}
\end{frame}

\begin{frame}{Estimators, validated by simulation}
  \begin{columns}[T]
    \column{0.42\textwidth}
    \small
    \begin{itemize}
      \item $\Hz$: mixed fBm on aggregated variances; control: sign autocorrelation.
      \item $H_{\mathrm{vol}}$: moments of the volume rate / autocovariance GMM.
      \item $H_\sigma$: GJR moments / Whittle on daily RV.
      \item $\delta$: same-side runs (proxy, exploratory).
    \end{itemize}
    Pre-registered rule: an estimator is used only where its simulated bias is below 0.02.
    \column{0.56\textwidth}
    \centering\fig{fig5_bias_maps.pdf}{0.72}
  \end{columns}
\end{frame}

\begin{frame}{$\Hz$, asset by asset}
  \centering\fig{fig1_h0_by_asset.pdf}{0.78}
\end{frame}

\begin{frame}{H1 volume $H_{\mathrm{vol}} = \Hz - \tfrac12$ \quad H2 volatility $H_\sigma = 2\Hz - \tfrac32$}
  \begin{columns}[T]
    \column{0.49\textwidth}
    \centering\fig{fig2_h1_hvol.pdf}{0.55}\par
    \small H1 (spot, confirmation): \alert{\resultHOne}
    \column{0.49\textwidth}
    \centering\fig{fig3_h2_hsigma.pdf}{0.55}\par
    \small H2 (spot, confirmation): \alert{\resultHTwo}
  \end{columns}
\end{frame}

\begin{frame}{H3 --- impact (exploratory), spot vs perpetual, robustness}
  \begin{columns}[T]
    \column{0.49\textwidth}
    \centering\fig{fig4_h3_delta.pdf}{0.55}\par
    \small Impact: \alert{\resultHThree}
    \column{0.49\textwidth}
    \centering\fig{fig8_specification_curve.pdf}{0.55}\par
    \small every variant, sorted, with its CI (not corrected)
  \end{columns}
\end{frame}

\begin{frame}{Does roughness forecast? (extension F)}
\begin{itemize}
  \item Out-of-sample forecasts of daily RV and volume, 1, 5, 22 days, 5-year history
  \item Exact rough predictor: Gaussian conditioning under day-averaged fBm + noise
  \item F1 rough vs HAR: \resultFOne \quad F2 order-flow $H$ vs volatility $H$: \resultFTwo
  \item Simulated rough volatility: $-8.5\,\%$ QLIKE vs HAR at 22 days; the published kernel does not beat HAR
\end{itemize}
\IfFileExists{figures/fig9_forecast_ratios.pdf}{\includegraphics[width=0.9\linewidth]{fig9_forecast_ratios.pdf}}{}
\end{frame}

\begin{frame}{Take-aways and reproducibility}
  \begin{itemize}
    \item Volume: \resultHOne; volatility: \resultHTwo; impact: \resultHThree.
    \item A rejection is reported with the same care as a confirmation; an undecided
      relation comes with the number of days needed to decide it.
    \item \texttt{make all} rebuilds every table and figure from a clean clone; the pipeline
      runs every day in a container and updates the rolling estimates and the dashboard.
    \item Code MIT; data and derived results CC BY-NC-SA 4.0 (Binance Vision).
  \end{itemize}
\end{frame}

\end{document}
